% Header: Here are all packages used and some additional definitions
%%%%%%%%%%%%%%%%%%%%%%%%%%%%%%%%%%%%%%%%%%%%%%%%%%%%%%%%%%%%%%%%%%%

\documentclass[11pt,a4paper]{article}
\usepackage[margin=2.5cm]{geometry}
\usepackage[onehalfspacing]{setspace}
\usepackage{graphicx} % zum Einbinden von Graphiken
\usepackage[breaklinks=true,colorlinks=true,linkcolor=blue,urlcolor=blue,citecolor=blue]{hyperref} % f. Referenzen
\usepackage{amsmath,amsthm,amssymb} % Mathematik Umgebung 
\DeclareMathOperator*{\argmax}{arg\,max}
\DeclareMathOperator*{\argmin}{arg\,min}


\usepackage{icomma} % Intelligentes Komma, das den richtigen Abstand zwischen Dezimalzahlen als auch in Formeln wählt.
%\usepackage[ngerman]{babel} % Deutsche Bezeichnungen bei Inhaltsangabe etc
\usepackage[T1]{fontenc}    % andere Schriftsatzkodierung für richtige Silbentrennung bei Umlauten
\usepackage[locale = DE,space-before-unit=true,per-mode = symbol]{siunitx} % Bessere Einheiten
\usepackage{placeins} % Definiert den Befehl “\FloatBarrier”, der die Ausgabe der davor eingebundenen Bilder erzwingt, befor der Text weiter geht. (Mit vorsicht zu verwenden)
%\usepackage[natbib,abbreviate=true,doi=false,style=numeric-comp,giveninits=true,sorting=none]{biblatex} % Modernes Paket zur Erzeugung von Bibliografien (benötigt biber!)
\usepackage[normalem]{ulem}
%\addbibresource{MyBibliography.bib} % Ort der .bib Datei, die die Datenbank für Literatur/Referenzen enthält.
\usepackage[ruled,vlined]{algorithm2e}
\SetKwInput{KwInput}{Input} % Set the Input
\SetKwInput{KwOutput}{Output} % set the Output
\graphicspath{{Bilder/}}

%\DeclareSIUnit{\dBm}{dBm}
%\DeclareSIUnit[per-mode=reciprocal]\WN{\per\centi\meter}


%Bibliography and citations
\usepackage[backend=biber, style=authoryear, sorting=nyt, maxnames = 10, maxcitenames = 2]{biblatex}
\addbibresource{references.bib}


\DeclareFieldFormat{citehyperref}{%
  \DeclareFieldAlias{bibhyperref}{noformat}%
  \bibhyperref{#1}}
\savebibmacro{cite}
\renewbibmacro*{cite}{%
  \printtext[citehyperref]{%
    \restorebibmacro{cite}%
    \usebibmacro{cite}}}

\let\cite\parencite

\AtBeginBibliography{\renewbibmacro*{begentry}{\printnames{author}}}
\AtEveryBibitem{\clearfield{note}}

\setlength{\bibhang}{1em}
%\defbibenvironment{bibliography}
%  {\list
%     {\printtext[labelnumberwidth]{%
%        \printfield{labelprefix}%
%        \printfield{labelnumber}}}
%     {\setlength{\labelwidth}{\labelnumberwidth}%
%     \setlength{\leftmargin}{\bibhang}
%     % \setlength{\leftmargin}{\labelwidth}%
%      \setlength{\labelsep}{\biblabelsep}%
%      \setlength{\itemsep}{\bibitemsep}%
%      \setlength{\parsep}{\bibparsep}%
%      \addtolength{\leftmargin}{\labelsep}%
%      \setlength{\itemindent}{-\bibhang}
%      %\setlength{\itemindent}{0pt}%
%      \renewcommand*{\makelabel}[1]{\hss##1}%
%      \raggedright}}
%  {\endlist}
%  {\item}

 

%%%%%%%%%%%%%%%%%%%%%%%%%%%%%%%%%%%%%%%%%%%%%%%%%%%%%%%%%%%%%%%%%%%
\begin{document}
%
\title{Bayesian Hierarchical CRR}
\author{Sophie Huebler \thanks{\href{mailto:sophie.huebler@utah.edu}{sophie.huebler@utah.edu}}}
\date{\today}
\maketitle
%\vfill
\pagenumbering{arabic} 
\tableofcontents

\newpage
\section{Thoughts}
\begin{itemize}
    \item Should the goal of the simulation scenario be based on the model itself, examining robustness and situations, or should it be matching the motivating problem?
    \item If showing the bounds, maintain variable selection as the main parameter of interest. Move towards incorporating correlation structure, different types of censoring, joint vs mutual predictor sets.
    \item If interested in the motivating problem, consider a copula model with known joint distribution. However, competing risks not exactly met. Come up with a better motivating scenario and find template paper?
    \item Proportional hazards assumption always broken in high dimension. 
\end{itemize}

\section{Separable Effects}
Background is competing risks. This is used to study the causal effect of a treatment on the event of interest. Separable is the direct effect on the treatment not mediated by the competing event.  \cite{Stensrud2019}.

\section{Subdistribution vs cause specific vs marginal}
\begin{itemize}
    \item \cite{Emura2020}
    \item Subhazard is always less than or equal to the marginal hazard.
    \item Non-identifiability of marginal distributions of two outcomes from observed data alone \cite{Tsiatis1975}
    \item Identifiable if assumed independence.
    \item Copula models bivariate survival distributions. Continuous bivariate on 0 to 1. Uniform marginals. 
\end{itemize}

\section{Example simulation}
\begin{itemize}
\item \cite{Binder2009}
    \begin{itemize}
    \item Fixed p, n, correlation structure, number of active covariates
    \item Informative covariates located in different correlation blocks
    \item Different informative covariates may or may not effect both events in either directin
    \item Event times generated from cause specific hazards (cox exponential)
    \item Censoring times from a uniform distribution
    \end{itemize}
\item \cite{EmuraChen2016}
\begin{itemize}
    \item Suggested alternatives: \cite{Binder2009, Chen2010}, random forest \cite{MogensenGerds2013}, univariate
    \item Copula model for dependent censoring (Frank or Clayton)
    \item Similar covariate structure to Binder
\end{itemize}
\end{itemize}

\section{Planning}
\begin{itemize}
    \item 
\end{itemize}



\newpage
\printbibliography[heading=bibintoc]





\end{document}

